\chapter{绪论}


\section{研究背景与意义}

\subsection{天地一体化网络的发展趋势与愿景}

随着信息技术的飞速发展和全球互联需求的日益增长，构建一个能够覆盖全球、提供无缝连接的网络体系已成为人类社会发展的战略性目标。传统的地面通信网络，尽管在人口密集区域实现了高度发达的覆盖，但其局限性也日益凸显：全球约有75\%的陆地面积和90\%以上的海洋区域仍是地面蜂窝网络的覆盖盲区。为了弥合这一“数字鸿沟”，实现真正意义上的“万物互联”，学术界与工业界将目光投向了广袤的太空，孕育出了天地一体化信息网络（Space-Terrestrial Integrated Network, STIN）这一宏伟构想。天地一体化信息网络并非简单的网络延伸，而是一个以天基网络为主体、地面网络为基础，能够支持陆、海、空、天各类用户随遇接入、按需服务的革命性信息基础设施。

这一愿景已从概念走向实践，并被提升至国家战略高度。中国在其航天白皮书中明确提出，要建设天地一体化信息网络，基本建成空间基础设施体系，并将其视为维护和拓展国家核心安全利益的关键举措。与此同时，全球范围内的商业航天公司与通信巨头也在积极布局。国际电信联盟（ITU）、欧洲空间局（ESA）等国际组织亦在推动相关技术研究与标准制定，旨在定义卫星通信与地面网络融合的业务交付架构，确保在多域供应环境中提供有保障的服务质量（QoS）与体验质量（QoE）。

天地一体化网络的发展与下一代移动通信技术，即第六代移动通信（6G）的演进紧密相连。如果说第五代移动通信（5G）通过其非地面网络（Non-Terrestrial Network, NTN）标准，初步探索了卫星等非地面网络与地面网络的协同工作模式，那么6G则致力于实现二者更深层次的无缝融合。在6G的蓝图中，卫星网络不再仅仅是地面网络的补充或备用链路，而是作为一个内生的、有机的组成部分，共同构成一个异构、立体、智能化的全球网络。这一融合网络旨在实现超高速率、超低延迟和超大连接数的目标，支持全息通信、数字孪生、全球物联网等前沿应用，最终达成“全球覆盖、随遇接入、按需服务、安全可信”的终极愿景。这种从“协同”到“融合”的范式转变，标志着天地一体化网络正从一个补充性技术方案，演进为未来全球信息社会不可或缺的核心基础设施。

\subsection{卫星物联网 (Satellite IoT) 数据激增带来的挑战}

天地一体化网络愿景的实现，极大地催生了卫星物联网（Satellite Internet of Things, SIoT）市场的蓬勃发展。卫星物联网利用卫星通信技术，为地面网络无法覆盖的偏远地区（如海洋、沙漠、山脉）和特殊行业（如农业、物流、海事、应急管理）的物联网设备提供连接服务，是实现全球物联网无缝覆盖的关键拼图。近年来，一场深刻的技术变革正在重塑卫星产业的成本结构，从而引爆了卫星物联网市场的指数级增长。这场变革的核心，是从传统上依赖于部署在地球静止轨道（Geostationary Orbit, GEO）的、重达数吨、造价高昂的大型卫星，转向在低地球轨道（Low Earth Orbit, LEO）部署由成百上千颗小型化、模块化、标准化生产的卫星组成的巨型星座。这一转变，加之可重复使用运载火箭技术带来的发射成本大幅下降，使得卫星连接的单位成本显著降低，从而将卫星物联网的应用门槛推向了前所未有的低点。

成本的下降直接驱动了市场的爆炸性扩张。根据多家市场研究机构的预测，全球卫星物联网市场正经历“爆炸式”增长阶段。IoT Analytics的报告显示，全球卫星物联网连接数在2024年已达到750万，并预测到2030年，相关市场规模将超过47亿美元，年复合增长率（CAGR）高达26\%。Berg Insight的预测更为激进，预计全球卫星物联网用户数将从2023年的510万飙升至2028年的2670万，复合年增长率接近40\%。这一增长趋势得到了业界的广泛共识，它预示着在不久的将来，将有数千万甚至上亿的物联网设备通过卫星接入网络，产生海量的传感、监控和遥测数据。

然而，这种由技术进步和成本下降驱动的良性循环，正催生出一个严峻的挑战：数据激增与星地链路（Satellite-to-Ground Link）带宽瓶颈之间的深刻矛盾。卫星物联网的成功部署，意味着全球范围内数以亿计的传感器将以前所未有的规模生成数据。这些数据，尽管单个数据包可能很小，但其汇聚而成的总量将是惊人的。据预测，到2025年，全球每天产生的数据量将达到463艾字节（Exabytes）。卫星物联网将是这一数据洪流的重要来源之一。然而，卫星与地面站之间的通信链路，受限于有限的频谱资源、高昂的建设成本以及物理定律，其总带宽是有限且宝贵的。当海量的原始物联网数据，包括大量冗余、低价值甚至错误的数据，都试图通过这一“窄门”回传至地面数据中心进行处理时，必然会导致严重的网络拥塞、传输延迟和服务质量下降。这种潜在的架构性危机，即网络的成功应用反过来造成了其自身基础设施的拥塞，正成为制约卫星物联网可持续发展的核心瓶颈。因此，如何高效地管理和优化这一海量数据流，在数据离开卫星之前进行有效的预处理，已成为该领域亟待解决的关键科学问题。

\subsection{LEO卫星网络资源受限的本质}

低地球轨道（LEO）卫星网络凭借其低延迟、低损耗的优势，成为构建下一代天地一体化网络和卫星物联网的核心技术。然而，与地面数据中心或网络设备相比，LEO卫星本质上是一个资源极度受限的平台。这种限制并非偶然，而是由空间环境的物理定律和航天工程的经济学原理共同决定的，其核心可以归结为对尺寸、重量和功率（Size, Weight, and Power, SWaP）的严苛约束。

发射成本是制约卫星设计的根本因素。将一公斤有效载荷送入LEO轨道仍需数千美元的成本，这迫使卫星设计者必须在功能与重量之间做出艰难的权衡，每一克重量都需精打细算。这种对重量的极致追求，直接限制了卫星上所有部件的规模和复杂性。其次，功率是卫星需要主要考量的因素。LEO卫星的主要能源来自太阳能帆板，其面积受限于卫星的整体尺寸（例如，立方星CubeSat的标准化尺寸）和发射整流罩的空间。有限的帆板面积意味着有限的发电功率，这一功率预算必须在通信载荷、计算单元、姿态控制系统和热管理系统等多个耗电大户之间进行分配。典型的LEO卫星最大发射功率可能仅为100瓦左右，远低于地面设备。

在这样严苛的SWaP约束下，卫星的各项资源都受到了显著限制。在计算资源方面，高性能的计算芯片通常功耗巨大，且对散热要求高。因此，星上处理器通常选用经过特殊设计、具备高能效比和抗辐射能力的宇航级CPU或FPGA。这些芯片的计算能力往往落后于地面商用芯片数代，其处理能力有限。例如，相关研究中的仿真参数显示，单颗LEO卫星的计算资源可能在3×10\textsuperscript{10} cycle/s的量级，显著低于地面云服务器的1×10\textsuperscript{11} cycle/s。在通信资源方面，频谱是稀缺且受到严格管制的资源。同时，天线的尺寸和发射功率直接影响通信带宽和信号质量，而这两者同样受到SWaP的严格限制。学术研究中，LEO网络的可用链路带宽通常被设定在几十兆赫兹的范围内，如25 MHz。此外，存储资源也同样受限，卫星无法像地面数据中心那样拥有海量的存储空间。

综上所述，LEO卫星的资源受限并非一个可以通过简单堆砌硬件来解决的工程问题，而是由发射经济学和轨道物理学决定的固有属性。这种计算、通信、存储和功率资源的全面受限，与前述卫星物联网的海量数据处理需求之间，形成了尖锐的矛盾。这一矛盾凸显了在卫星上采用轻量级、高效率的软件和算法方案，以“四两拨千斤”的方式提升系统整体性能的极端重要性和紧迫性。

\subsection{在网计算赋能星上边缘处理的研究意义}

面对卫星物联网数据激增与LEO卫星资源受限的双重挑战，传统的卫星网络架构显得力不从心。传统的“透明转发”（Bent-pipe）架构将卫星视为一个简单的空中中继器，其唯一功能是将地面终端的信号接收、放大并转发至地面站，所有的数据处理和智能分析都在地面数据中心完成。这种模式虽然技术成熟、卫星设计简单，但在数据量爆炸的今天，它将海量的原始数据，包括其中大量的冗余、错误和低价值信息，不加区分地挤占宝贵的星地链路带宽，造成了巨大的资源浪费和网络拥塞。

为了打破这一困局，将计算能力推向网络边缘，在数据源头进行处理的边缘计算（Edge Computing）思想被引入到卫星网络领域，催生了“星上边缘计算”（Satellite Edge Computing）或“轨道边缘计算”（Orbital Edge Computing）这一新兴范式。其核心思想是将卫星从一个“哑管道”转变为一个智能的“边缘节点”，赋予其在轨进行数据处理、分析和决策的能力。在这一范式下，一种更具革命性的技术——在网计算（In-Network Computing, INC），展现出巨大的潜力。在网计算利用可编程网络设备（如P4交换机）的数据平面，在数据包转发的路径上，以线速（line-rate）执行计算任务，实现了计算与网络的深度融合。

将在网计算技术应用于星上边缘处理，意味着我们可以将卫星内部的交换芯片从一个单纯的“路由器”升级为一个高效的“数据预处理器”。其研究意义重大且深远。首先，最直接的意义在于极大地缓解星地链路的带宽压力。通过在卫星上部署轻量级的数据去重、压缩、过滤等在网计算任务，可以在数据回传前剔除大量冗余信息，只将高价值、经过提炼的结果传回地面。中国的“天算星座”计划已经通过星地协同计算的初步验证，实现了回传数据量减少90\%的惊人效果，这充分证明了该路径的巨大潜力。其次，在网计算能够显著降低业务延迟。对于海洋监测、灾害预警等时间敏感型应用，将数据分析任务从地面数据中心前移至卫星，可以免去长距离传输的往返时延，实现近乎实时的响应。再者，它能有效提升数据安全与隐私保护水平。敏感数据无需在开放的地面网络中传输，在轨处理后即可丢弃原始数据，大大降低了数据泄露的风险。

最终，在网计算赋能的星上边缘处理，是从根本上解决卫星物联网规模化发展瓶颈的关键技术路径。它通过将计算负载智能地分布到网络边缘，优化了整个系统的资源利用效率，有望使构建一个经济、高效、智能、安全的全球卫星物联网成为可能。因此，研究面向LEO卫星物联网场景、基于在网计算的星上数据预处理机制，不仅具有重要的理论价值，更对推动天地一体化网络和6G技术的落地应用具有不可估量的实践意义。

\section{国内外研究现状}

\subsection{LEO卫星网络技术研究现状}

近年来，随着SpaceX的Starlink、OneWeb、Amazon的Kuiper等巨型星座计划的推进，LEO卫星网络技术的研究进入了一个新的黄金时代，呈现出理论与实践并进的繁荣景象。当前的研究现状主要围绕网络架构、标准化进程、以及关键技术挑战展开。

在网络架构方面，现代LEO卫星网络正从早期的单点通信模式，向着构建一个具有强大在轨处理与路由能力的太空互联网（Internet in Space）演进。其核心特征是大规模采用星间链路（Inter-Satellite Links, ISLs），特别是高速激光星间链路。通过ISL，轨道上的卫星可以构建成一个动态的、高带宽的网状网络（Mesh Network），数据可以在卫星之间进行多跳转发，从而实现全球范围内的低延迟通信，减少对地面站的依赖。这种架构使得网络流量可以“绕行”拥堵的地面网络，为偏远地区间的通信提供更优的路径。Starlink的实践表明，其网络会以15秒为周期进行全局性的网络重构，动态调整用户与卫星、卫星与地面站之间的连接路径，以适应卫星的快速移动和实现负载均衡，这种高度动态性是当前LEO网络架构研究的核心特征之一。

在标准化进程方面，为了实现卫星网络与地面移动通信网络的无缝融合，各大标准组织，特别是第三代合作伙伴计划（3GPP），正在积极推进非地面网络（NTN）的标准化工作。自Release-15开始进行场景和信道模型研究以来，3GPP在Release-17中正式完成了支持NTN的首个规范版本，标志着卫星通信被正式纳入5G生态系统。Release-17主要聚焦于采用透明转发（Transparent Payload）模式的卫星，即卫星仅作为射频信号的中继，而基站（gNB）功能仍在地面，这种方式对现有5G NR协议的改动最小。其研究内容涵盖了为适应卫星通信长时延、大频偏等问题而进行的物理层和高层协议增强，以及移动性管理、服务质量（QoS）保障等关键系统问题。随后的Release-18及未来的6G研究，将进一步探索更复杂的场景，如具有在轨再生处理能力（Regenerative Payload）的卫星、星上基站、以及更高效的移动性和切换管理机制，旨在让卫星接入成为与地面Wi-Fi、蜂窝网络同等地位的一种标准接入方式。

在关键技术挑战的研究上，学术界正从多个维度展开探索。动态路由是LEO网络研究的核心难点之一，由于拓扑的周期性剧烈变化，传统的地面网络路由协议（如OSPF）难以直接适用，研究人员正在探索基于时间演化图、或利用拓扑可预测性进行改进的路由算法。移动性管理是另一大挑战，包括小区级的移动性（波束切换）和卫星级的移动性（卫星切换），如何实现无缝、低延迟的切换是保障用户体验的关键。网络安全问题也日益受到重视，LEO网络广阔的覆盖范围和广播特性使其容易成为DDoS攻击等网络攻击的目标，研究如何在这种动态、资源受限的环境下设计高效的安全防护机制是一个活跃的研究方向。此外，由卫星高速运动引起的多普勒频移补偿、信道建模、以及面向物联网场景的随机接入协议优化等问题，也是当前LEO网络技术研究的热点。

\subsection{在网计算与P4可编程网络技术研究现状}

在网计算（In-Network Computing, INC）作为一种新兴的计算范式，旨在将计算任务从终端设备或中心云，卸载到网络路径中的交换机、路由器等转发设备上执行，其发展与数据平面可编程技术的演进密不可分。这一领域的研究现状，标志着网络从一个被动的数据管道，向一个主动的、可编程的计算平台转变。

研究的起点可以追溯到软件定义网络（Software-Defined Networking, SDN）的出现。SDN通过解耦控制平面与数据平面，首次实现了对网络行为的集中式可编程控制。其代表性协议OpenFlow允许控制平面动态下发流表规则，指导数据平面进行“匹配-动作”（Match-Action）处理。然而，OpenFlow的局限性也十分明显：它所能匹配的协议头字段集合是固定的，由标准预先定义。当需要处理新的协议或自定义协议头时，必须等待标准更新和硬件厂商支持，这极大地束缚了数据平面的创新速度和灵活性。

为了突破OpenFlow的桎梏，P4（Programming Protocol-independent Packet Processors）语言应运而生，并迅速成为数据平面可编程领域的事实标准。P4的革命性在于它实现了三个核心目标：协议无关性、现场可重构性和目标可移植性。协议无关性意味着网络程序员可以自行定义数据包的解析逻辑，让交换机能够识别和处理任意格式的协议头，而不再受限于厂商的硬编码实现。这种“自顶向下”的设计范式，将数据平面的定义权从硬件厂商交还给了网络运营者和研究者，极大地加速了网络功能的创新周期。

基于P4的在网计算已经在地面网络，特别是数据中心网络中，展现出巨大的应用价值。当前的研究主要集中在以下几个方面：首先是网络遥测与监控。通过P4实现的带内网络遥测（In-band Network Telemetry, INT）技术，可以将交换机内部的状态信息（如队列深度、处理时延）直接编码到数据包中，随包传输，从而实现对网络状态的线速、全路径、高精度可见性。其次是流量工程与优化。利用P4的灵活性，研究人员设计并实现了多种高效的负载均衡、拥塞控制和主动队列管理（AQM）算法，这些算法可以直接在数据平面运行，实现微秒级的快速响应。再次是网络安全。P4交换机被用于构建高性能的DDoS攻击缓解系统和可编程防火墙，它们能够在网络入口处以线速检测并过滤恶意流量，远比传统的基于CPU的方案高效。此外，一些更前沿的研究开始探索利用P4数据平面来加速分布式系统应用，如分布式共识协议、键值存储和机器学习推理等。

为了支撑P4的研发与应用，一个完整的生态系统已经形成，包括了官方的P4语言规范、开源的编译器后端（如p4c）、行为级软件交换机模拟器（如BMv2）、以及与网络仿真平台（如Mininet）的集成工具，这些都极大地降低了研究和开发P4应用的门槛。总而言之，P4和在网计算技术已经从一个学术概念发展成为一个拥有坚实理论基础、丰富应用场景和完善工具链的成熟研究领域。

\subsection{物联网数据聚合与预处理技术研究现状}

物联网（IoT）场景下，海量设备产生的原始数据往往具有高度的冗余性、时空相关性和价值稀疏性。直接传输和存储所有原始数据会造成巨大的通信带宽和存储资源浪费。因此，在数据传输路径上对数据进行聚合与预处理，是物联网系统设计的关键环节。当前，该领域的研究主要聚焦于设计轻量级的算法，以适应物联网终端和边缘节点资源受限的特点。

数据聚合技术旨在通过融合来自多个传感器的数据来消除冗余传输，从而节约能源和带宽。根据实现方式，聚合技术可分为多种。一种是基于树状或簇状的层次化聚合，网络中的节点被组织成簇，簇内数据先汇聚到簇头节点，再由簇头节点进行聚合后向上传输，这种方式可以有效减少传输跳数和数据总量。另一种是数据中心路由（Data-centric Routing），其中数据本身的内容（如事件类型）决定了其路由路径，典型的例子是“定向扩散”（Directed Diffusion），汇聚节点（Sink）通过发布兴趣来“拉取”相关数据，中间节点可以对数据进行缓存和融合。这些聚合方法的核心目标都是在保证信息质量的前提下，最大化网络生命周期。

数据预处理技术则更侧重于对数据内容本身进行清洗和精炼。其中，数据去重（Deduplication）是研究的热点之一。由于物联网设备部署密集，相邻设备在同一时刻采集到的数据往往高度相似或完全相同。去重技术通过识别并消除这些重复数据来节省资源。现有技术主要包括：基于哈希的去重，通过计算数据块的哈希值（如SHA-1）来判断其是否重复，这是最常用的方法。为了处理相似但不完全相同的数据，研究人员提出了相似哈希（SimHash）或基于内容分块（Content-Defined Chunking）结合布隆过滤器（Bloom Filter）的方案，以识别相似数据。此外，为了在去重的同时保护数据隐私，基于同态加密或安全多方计算的加密去重方案也被提出，但其计算开销较大。一些研究还利用默克尔树（Merkle Tree）或其变体（如扩展默克尔网格）来高效地进行数据校验和去重验证。

除了去重，数据过滤也是重要的预处理手段。这包括滤除明显的异常值或由传感器故障产生的错误数据。近年来，随着人工智能技术的发展，轻量级的机器学习模型被用于物联网边缘节点，进行更智能的数据预处理。例如，利用决策树、支持向量机（SVM）或轻量级的深度学习模型（如1D-CNN）在边缘端进行数据分类或异常检测，只将有意义的事件或异常模式上传，从而大幅减少数据量。联邦学习（Federated Learning）作为一种分布式机器学习范式，也为物联网数据处理提供了新思路，它允许在不上传原始数据的情况下，在各设备上进行本地模型训练，仅聚合模型参数，兼顾了效率与隐私。

总的来看，物联网数据聚合与预处理技术的研究已经相当深入，形成了丰富的算法库。然而，这些研究的一个共同特点是，其算法实现大多是面向通用的计算平台，如物联网设备内置的微控制器（MCU）或边缘服务器的CPU。它们的设计假设了一个可以执行复杂逻辑、循环、分支和拥有灵活内存访问的冯·诺依曼计算模型。

\subsection{现状评述与研究切入点分析}

通过对LEO卫星网络、在网计算与P4技术、以及物联网数据预处理技术这三个相关领域研究现状的梳理，可以清晰地看到一幅技术发展的宏观图景：一方面，LEO卫星网络正在成为承载全球物联网的关键基础设施，但其面临着由数据激增和自身资源受限所导致的严峻挑战；另一方面，P4可编程网络和在网计算技术为在网络内部高效处理数据提供了前所未有的强大工具，并在地面网络中取得了巨大成功；同时，针对物联网数据的轻量级聚合与预处理算法也已发展得相当成熟。然而，当我们将这三个领域进行交叉审视时，一个明显的研究空白与机遇便浮现出来。

现有研究呈现出一种“割裂”的状态。关于LEO网络的研究，大多集中在物理层、链路层和网络层的协议设计与优化上，如路由、移动性管理和信道接入，较少深入探讨如何利用网络层以上的能力对承载的数据内容进行智能处理。关于P4在网计算的研究，其应用场景几乎全部局限于高带宽、资源相对充裕、拓扑稳定的地面数据中心或企业网，鲜有工作将其置于LEO卫星网络这种高动态、资源极度受限的极端环境中进行考量。而关于物联网数据预处理算法的研究，其设计和实现完全基于传统的CPU计算模型，未曾考虑过如何将其适配到P4所代表的、高度受限的流水线式数据平面编程模型中。

正是这三个领域的交叉点，构成了本论文的核心研究切入点。将P4在网计算技术应用于LEO卫星网络，以解决卫星物联网的数据预处理问题，是一个极具前瞻性但又充满挑战的课题。现有技术无法直接迁移，必须进行根本性的创新。具体而言，本研究旨在填补以下关键空白：

首先，\textbf{算法层面的创新}。现有的物联网去重、过滤算法，无论是基于复杂的密码学操作，还是基于迭代式的机器学习模型，都无法直接在P4数据平面中实现。P4的流水线架构不支持循环，状态存储（寄存器、计数器）资源极为有限，且可执行的算术逻辑单元（ALU）操作非常基础。因此，本研究的第一个切入点是：如何重新设计数据预处理算法，使其逻辑能够被“翻译”成P4的“匹配-动作”表（Match-Action Table）流水线？这要求算法必须是“单遍”（single-pass）的，并且对状态存储和计算复杂度的要求极低。

其次，\textbf{系统层面的挑战}。即便设计出P4兼容的算法，如何在一个高速运动、拓扑每时每刻都在变化的LEO网络中有效管理其所需的状态，是另一个巨大的挑战。例如，一个在网去重应用需要在P4交换机中维护一个已见数据包哈希值的表。当卫星网络发生切换或重构时（如Starlink的15秒周期性重构），这个状态表是应该被丢弃、迁移还是与其他卫星同步？状态的生命周期管理、一致性保证以及在动态网络中的高效同步，是现有P4应用从未遇到过的问题。这是本研究的第二个切入点。

综上所述，本论文的研究切入点在于探索上述三个领域的深度融合，致力于回答一个核心问题：如何设计一个面向LEO卫星物联网场景的、基于P4在网计算的轻量级数据预处理框架，并解决其中的关键算法设计与动态状态管理难题。这项研究不仅有望为解决卫星物联网的带宽瓶颈问题提供一个全新的、高效的解决方案，也将极大地拓展P4在网计算的应用边界，推动其从地面网络走向广阔的空天领域。

\section{本文主要研究内容}

基于上述研究背景、意义及切入点分析，本文将围绕在LEO卫星网络中利用在网计算技术对物联网数据进行高效预处理这一核心目标，展开以下两个层面的主要研究内容。

\subsection{星上轻量级数据预处理机制设计}

本部分内容聚焦于算法与机制的顶层设计，旨在解决传统物联网数据处理算法与P4数据平面编程模型不兼容的核心矛盾。研究将从卫星物联网数据的典型特征（如数据冗余、格式多样性）出发，探索能够在P4交换机严苛的资源限制下高效执行的数据预处理操作。主要研究内容包括：

第一，\textbf{在网数据去重算法设计}。针对物联网传感器数据中普遍存在的大量重复信息，研究将设计一种或多种轻量级的在网去重算法。这需要摒弃传统软件实现中复杂的逻辑，将去重过程分解为P4流水线能够执行的基本操作。研究将探索如何利用P4的哈希计算、寄存器阵列和布隆过滤器等原生能力，构建一个内存占用小、匹配速度快的去重状态表，并设计相应的匹配-动作逻辑，以线速识别并丢弃重复的数据包。

第二，\textbf{在网数据过滤机制研究}。除了完全重复的数据，物联网数据流中还包含大量根据特定规则可以被过滤的“无用”数据（例如，超出正常阈值的异常读数、不符合特定格式的数据包等）。本研究将设计一种灵活的在网过滤机制，允许网络管理员通过控制平面动态配置过滤规则。这些规则将被编译成P4的匹配表项，使得卫星上的P4交换机能够根据数据包的头部或载荷内容，实时地对数据流进行筛选。

第三，\textbf{面向异构物联网数据的可编程解析}。物联网协议种类繁多，标准不一。本研究将利用P4的协议无关性，设计一个可编程的解析器，使其能够根据数据包的特征动态识别并解析不同类型的物联网数据格式，从而为后续的去重和过滤操作提取出关键的元数据字段。

\subsection{面向异构任务的星上计算卸载与资源调度策略}

在设计了具体的预处理算法之后，本部分内容将从系统层面出发，研究如何在一个动态的、多任务的星上环境中，对这些在网计算任务进行有效的资源分配和管理。这部分研究旨在解决星上计算资源稀缺性和网络拓扑高动态性带来的挑战。主要研究内容包括：

第一，\textbf{P4交换机资源感知与建模}。为了实现高效的资源调度，首先需要对P4交换机的内部资源（如匹配表项、寄存器内存、ALU单元等）进行精确的感知和建模。研究将分析P4程序编译后对硬件资源的占用情况，建立一个能够量化不同在网计算任务资源消耗的模型，为调度决策提供依据。

第二，\textbf{动态网络环境下的状态管理策略}。这是本研究的核心难点之一。针对LEO网络拓扑的周期性变化和频繁切换，研究将探索不同的在网状态管理策略。例如，对于去重应用所需的状态哈希表，是采用“无状态”设计（切换后状态即失效），还是设计一种轻量级的状态同步或迁移机制，在卫星切换时将关键状态信息传递给下一个服务卫星。研究将分析不同策略在开销、一致性和处理效率之间的权衡。

第三，\textbf{异构任务联合调度与卸载决策}。星上平台可能需要同时承载多种数据预处理任务（如去重、过滤、聚合等），以及其他网络功能。这些任务对P4资源的竞争关系需要被妥善管理。本研究将探索一种联合调度策略，根据任务的优先级、数据流的特征以及当前卫星的可用资源，动态地决定哪些任务被加载到P4流水线中，以及如何为它们分配合适比例的硬件资源，以实现系统整体效用的最大化。

\section{本文主要贡献与创新}


\section{论文结构安排}


